\section{三重积、Q 值与成本缩放}

上一章解释了“用磁场装住太阳”，本章进入判断聚变装置性能的核心指标。访谈中杨钊反复提到三重积和 Q 值。三重积是等离子体密度、温度和能量约束时间的乘积；Q 值是聚变输出功率除以输入功率。通俗说，三重积决定“这锅等离子体有没有足够好”，Q 值决定“这台装置是否有能量增益”。

\[
\mathcal{T}=n T \tau_E
\]

其中，$n$ 表示等离子体粒子密度，$T$ 表示温度，$\tau_E$ 表示能量约束时间，$\mathcal{T}$ 是三重积。访谈中提到，对于氘氚聚变，三重积达到约 $10^{21}$ 量级时接近 Q 等于 1 的区间，再往上 Q 会快速提升；若要做电站，工程上通常希望 Q 显著大于 1，例如 Q 等于 10 附近。

\[
Q=\frac{P_{\mathrm{fusion}}}{P_{\mathrm{input}}}
\]

其中，$P_{\mathrm{fusion}}$ 表示聚变输出功率，$P_{\mathrm{input}}$ 表示外部输入功率。Q 等于 1 是能量盈亏平衡的物理门槛，但发电站还要考虑取热、发电、循环和辅助系统损耗，因此实际设计目标要高得多。

\begin{importantbox}{Q 大于 1 还不等于电站盈利}
Q 等于 1 说明聚变输出功率等于输入功率，但电站需要把热转电、维持磁体和冷却系统、处理运行损耗。真正有工程意义的目标通常是 Q 显著大于 1，并且整站度电成本可接受。
\end{importantbox}

\subsection{为什么磁场如此关键}

本节把公式翻译成工程直觉。上一小节给出了三重积和 Q 值，但创业公司真正要优化的不是黑板上的一个符号，而是一台机器的半径、磁体、结构件、低温系统和建造成本。磁场之所以关键，是因为它同时牵动等离子体约束和装置尺寸，最后会传导到资本开支和度电成本。

三重积本身由密度、温度和约束时间构成，但工程上不能随便单独提高其中一个量，因为提高一个参数可能让另外两个参数变差。经验定标律告诉工程师，装置大小和磁场强度是更直接的杠杆。访谈中给出的直观关系是：托卡马克三重积大致正比于磁场的 3.5 次方，正比于装置线性尺寸的 2.5 次方。

\[
\mathcal{T}\propto B^{3.5}R^{2.5}
\]

其中，$B$ 表示磁场强度，$R$ 表示装置线性尺寸或大半径的量级。这个关系不是第一性原理公式，而是大量装置实验数据拟合出的定标律。它的商业含义很强：如果磁场能提高一倍，装置线性尺寸可以显著缩小；体积和造价又大致随尺寸三次方变化，因此高磁场可能带来数量级成本变化。

\begin{knowledgebox}{读图：为什么高磁场能换来低成本}
把磁场提高，看起来只是物理参数变化；但在工程系统里，它会连锁影响装置半径、磁体规模、真空室体积、结构材料、建造周期和资本开支。这就是高温超导被视作拐点材料的原因。
\end{knowledgebox}

\begin{knowledgebox}{从公式到成本：三步传导}
\begin{tabularx}{\textwidth}{p{0.20\textwidth}p{0.34\textwidth}X}
\toprule
层级 & 发生了什么 & 为什么重要 \\
\midrule
物理层 & 更高磁场提高约束能力 & 三重积更容易达到 Q 大于 1 或 Q 大于 10。 \\
装置层 & 同等性能下线性尺寸可以缩小 & 真空室、结构件、低温系统和安装空间都跟着缩小。 \\
商业层 & 体积和质量下降带来造价下降 & 度电成本才可能从科学装置水平进入电站水平。 \\
\bottomrule
\end{tabularx}
\end{knowledgebox}

\subsection{惯性约束为什么不直接变成民用电站}

访谈中也解释了惯性约束路线。美国 NIF 曾在靶丸层面实现 Q 大于 1 的结果，但这不是端到端电站意义的 Q：输入端通常指激光能量，而把电能转成激光本身有巨大损耗；同时，惯性约束反应是极短脉冲，不是稳态连续运行。因此它在武器物理和高能量密度研究中很重要，但并不是当前民用聚变发电的主流路线。

\begin{warningbox}{不要只看单个 Q 数字}
不同路线的 Q 值口径不同。惯性约束的靶丸增益、托卡马克等离子体增益、整站电力增益不是同一个概念。比较技术路线时，必须问清楚输入功率算到哪一层、输出能量是否可连续取出、系统能否长时间运行。
\end{warningbox}

\subsection{本章小结}

三重积和 Q 值是理解聚变性能的入口；磁场强度和装置尺寸是影响三重积的关键工程杠杆；高温超导的意义在于把“达到 Q 大于 10”的装置从超大科学工程压缩成更小、更快、更低成本的路径。

